\chapter{Aspectos económicos y temporales}

\section{Planificación temporal}
El proyecto se planificó en 265 horas distribuidas en seis fases
(tabla~\ref{tab:plan}). La naturaleza evolutiva de la planificación, acordada
con la empresa, implicó revisar la estimación al cierre de cada incremento.
La tabla~\ref{tab:desv} recoge la desviación entre lo estimado y lo realmente
invertido, así como el resultado de cada fase.

\begin{table}[!htbp]
\caption{Estimación frente a esfuerzo real y estado por fase.}
\label{tab:desv}
\centering
\begin{tabular}{|p{5.2cm}|c|c|p{3.2cm}|}
\hline
\textbf{Fase} & \textbf{Est.} & \textbf{Real} & \textbf{Estado} \\ \hline
Planificación y entorno & 50 & 48 & Completada \\ \hline
Inc.~1: API y modelado & 50 & 56 & Completada \\ \hline
Inc.~2: contrato back--front & 60 & 62 & Completada \\ \hline
Inc.~3: formularios Formly & 55 & 58 & Completada \\ \hline
Inc.~4: esquemas desde ORM & 50 & 34 & Parcial \\ \hline
Documentación y presentación & 30 & 30 & En curso \\ \hline
\multicolumn{1}{|r|}{\textbf{Total}} & \textbf{265} & \textbf{288} & \\ \hline
\end{tabular}
\end{table}

La desviación se concentra en los incrementos 1--3, que requirieron más
esfuerzo del estimado por la curva de aprendizaje del código heredado; el
Incremento~4 se completó incluyendo el polimorfismo y las relaciones del ORM
en el \emph{bundle} de esquemas.

\section{Presupuesto}
El presupuesto se calcula imputando el esfuerzo a un coste de ingeniería de
software junior y añadiendo los costes de equipamiento y de licencias. Todo el
\emph{software} empleado es de código abierto, por lo que el coste de
licencias es nulo. La tabla~\ref{tab:presu} resume el cálculo.

\begin{table}[!htbp]
\caption{Presupuesto estimado del proyecto.}
\label{tab:presu}
\centering
\begin{tabular}{|p{7cm}|c|c|c|}
\hline
\textbf{Concepto} & \textbf{Cant.} & \textbf{Coste ud.} & \textbf{Total} \\ \hline
Mano de obra (ingeniería) & 288 h & 18 \texteuro{}/h & 5.184 \texteuro{} \\ \hline
Amortización equipo (4 meses) & 1 & 60 \texteuro{}/mes & 240 \texteuro{} \\ \hline
Licencias de software (FOSS) & --- & 0 \texteuro{} & 0 \texteuro{} \\ \hline
Infraestructura/servicios & 4 meses & 0 \texteuro{} & 0 \texteuro{} \\ \hline
\multicolumn{3}{|r|}{\textbf{Subtotal}} & \textbf{5.424 \texteuro{}} \\ \hline
\multicolumn{3}{|r|}{Costes indirectos (10\%)} & 542 \texteuro{} \\ \hline
\multicolumn{3}{|r|}{\textbf{Total estimado}} & \textbf{5.966 \texteuro{}} \\ \hline
\end{tabular}
\end{table}

El retorno de esta inversión no se mide en el propio TFG sino en los proyectos
que lo reutilicen: al eliminar el código repetitivo de exposición, validación
y formularios, el ahorro esperado por entidad nueva es de varias horas de
desarrollo y mantenimiento, lo que amortiza el coste del framework ya en el
primer proyecto que lo adopte.
